% ---
% Capitulo de Introducao
% ---

\chapter{Introdução\label{cha:intro}}

A análise de dados é o processo de examinar um conjunto de observações com o objetivo de descrever a estrutura que o organiza, extrair dele informação que não está explícita nos registros individuais, e sustentar decisões a partir dessa informação. Ela está presente em atividades que produzem registros em volume, da detecção de intrusões em redes \cite{network-intrusion-detection} à identificação de anomalias em processos de manufatura \cite{manufacturing-anomaly-detection}, passando pela separação de objetos astronômicos em levantamentos fotométricos \cite{star-galaxy-separation}, pela previsão de geração eólica \cite{wind-power-prediction} e pelo monitoramento de processos industriais de soldagem \cite{welding-curve-discovery}. Sua importância cresce com o volume de dados: quanto maior o volume registrado, menor a fração que pode ser inspecionada manualmente, e maior a dependência de métodos que descrevam a estrutura do conjunto sem que ela precise ser conhecida de antemão.\\

Dentro da análise de dados, o agrupamento é a tarefa de particionar as observações em conjuntos coesos, e os métodos baseados em densidade formam uma família adequada à análise exploratória. Eles definem um grupo como uma região do espaço em que as observações se concentram acima de um limiar de densidade, e tratam como ruído aquilo que não atinge esse limiar. Essa definição confere três propriedades úteis. Primeiro, o número de grupos não precisa ser informado, pois decorre da própria distribuição. Segundo, os grupos não são restritos a formas convexas ou esféricas, o que permite capturar estruturas alongadas, aninhadas ou irregulares. Terceiro, o ruído é definido explicitamente, evitando se tornar um resíduo forçado a pertencer a algum grupo. Essas propriedades tornam os métodos baseados em densidade aplicáveis a tarefas que vão além do agrupamento propriamente dito, como discriminação de ruído, classificação, regressão e modelagem \cite{framework-density-clustering}.\\

O representante mais conhecido dessa família é o \ac{DBSCAN} \cite{dbscan}, que qualifica cada objeto com base na quantidade de vizinhos dentro de um raio $\varepsilon$ e conecta objetos densamente alcançáveis entre si. Seu sucessor hierárquico, o \ac{HDBSCAN*} \cite{hdbscan}, elimina a necessidade de fixar $\varepsilon$ ao construir a hierarquia completa dos agrupamentos obtidos para todos os valores possíveis desse limiar, deixando apenas um parâmetro livre, o $\mpts$, que define quantos vizinhos são considerados na estimativa de densidade local. Foi demonstrado que toda essa família pode ser descrita como instâncias de um mesmo arcabouço, cuja operação central é o cálculo da árvore geradora mínima (\ac{MST}) de um grafo completo em um espaço transformado \cite{framework-density-clustering}, em que o peso de cada aresta é uma estimativa de quão densamente conectados estão os dois objetos que ela liga. A hierarquia do \ac{HDBSCAN*}, os agrupamentos do \ac{DBSCAN} e as demais tarefas derivadas são todos extraíveis dessa única árvore.\\

Essa constatação expõe, contudo, um custo prático. O $\mpts$ não pode ser escolhido de forma principiada antes de olhar para os dados, e nenhum valor único é apropriado para conjuntos cujos grupos diferem em densidade, o que é comum. É usual, portanto, executar o método várias vezes e comparar as hierarquias resultantes. Como cada valor de $\mpts$ define um espaço transformado diferente, cada execução exige uma nova \ac{MST} sobre um grafo completo, ao custo de $O(n^2 d)$ para $n$ objetos de dimensionalidade $d$. O \ac{CORE-SG} \cite{core-sg} remove essa barreira. Ele é um grafo único, construído uma só vez, formado pela união do grafo de $k$ vizinhos mais próximos, o \ac{k-NNG}, com uma \ac{MST} calculada no maior valor de interesse do parâmetro, denotado $\kmax$. Seus autores demonstram que esse grafo contém uma \ac{MST} do grafo completo para todo $\mpts \leq \kmax$, de modo que extrair o resultado para um valor específico do parâmetro se reduz a repesar $O(n\kmax)$ arestas já armazenadas e executar um algoritmo de árvore geradora mínima sobre elas. Uma varredura inteira sobre o parâmetro passa a custar aproximadamente o mesmo que uma única execução do método original.\\

% ---
\section{O problema\label{sec:problem-statement}}
% ---

O \ac{CORE-SG} resolve o custo de repetir a análise, mas não altera o custo de construir a estrutura uma primeira vez. A limitação que este trabalho ataca está nessa construção.\\

A construção original segue um modelo de execução centralizado: as duas metades são calculadas sobre o conjunto inteiro, de uma só vez, em um único espaço de memória. A metade de árvore geradora agrava o quadro, por ser definida sobre um grafo completo de $O(n^2)$ pares. O maior conjunto processável fica assim determinado pela memória de uma única máquina, independentemente de quantas estejam disponíveis.\\

A limitação é mais severa em aceleradores gráficos, e a Seção \ref{sec:problem-scale} a desenvolve. A operação dominante na construção do \ac{k-NNG} é o cálculo de distâncias entre pares, aritmeticamente denso e independente, que é o perfil de carga em que uma \ac{GPU} rende mais. Como aceleradores típicos oferecem uma ordem de grandeza menos memória que a principal de um servidor, e a estrutura intermediária cresce com $n$ e com $\kmax$, a memória do dispositivo é o recurso que primeiro se esgota.\\

A consequência é um limite abrupto, situado em torno de $9{,}1 \times 10^6$ pontos para $d = 128$ e $\kmax = 32$ em um acelerador de 8 GB. Acrescentar aceleradores não desloca esse limite, porque uma construção monolítica não se divide entre dispositivos, de modo que o segundo permanece ocioso sem contribuir com sua memória.\\

O problema atacado é o seguinte: a construção do \ac{CORE-SG} é limitada pela memória de um único dispositivo, e acrescentar dispositivos não contorna essa limitação. O objetivo é construir, com os aceleradores disponíveis, grafos que excedem a memória de qualquer um deles isoladamente; a velocidade é secundária.\\

% ---
\section{Questões de pesquisa\label{sec:research-questions}}
% ---

O problema acima se desdobra em quatro questões, que organizam o restante do texto.\\

QP1. É possível decompor a construção do \ac{CORE-SG} em subproblemas cuja memória individual seja controlável de forma independente do tamanho do conjunto de dados? Uma decomposição útil precisa satisfazer duas condições simultâneas. A memória que cada subproblema exige deve ser função de um parâmetro sob controle do usuário, independentemente de $n$; e os subproblemas devem ser autocontidos o bastante para serem despachados a qualquer acelerador disponível sem comunicação entre dispositivos durante a execução, pois essa comunicação reintroduziria o acoplamento que a decomposição elimina.\\

QP2. Qual é o custo dessa decomposição em qualidade, e esse custo é controlável? Particionar um conjunto de dados para construir um grafo de vizinhança significa deixar de comparar todo par cujos dois pontos caiam em partes distintas. O grafo resultante é, portanto, uma aproximação pior do que aquela produzida por uma construção não particionada sobre os mesmos dados. É necessário quantificar essa perda, identificar o mecanismo que a recupera, e determinar se existe um parâmetro que permita ao usuário escolher sua posição na curva entre memória e qualidade.\\

QP3. A segunda metade do \ac{CORE-SG}, a árvore geradora mínima, pode ser substituída por uma estimativa cujo custo não dependa de $n$? Mesmo que o \ac{k-NNG} seja decomposto com sucesso, a \ac{MST} sobre o grafo completo permanece como um componente centralizado e quadrático. Se a estrutura inteira deve escalar, essa metade precisa ou ser distribuída, ou ser substituída por algo cujo custo seja governado por outra grandeza que não o número de pontos.\\

QP4. O que a decomposição custa em tempo, e o que ela oferece em contrapartida? Uma construção decomposta executa trabalho que uma construção monolítica não executa: a partição, a mesclagem dos resultados parciais e a recuperação dos pares não comparados. É preciso medir esse sobrecusto, determinar que fração da construção escala com a quantidade de aceleradores, e situar o resultado em relação às alternativas disponíveis.\\

% ---
\section{Abordagem e contribuições\label{sec:contributions}}
% ---

Os dois grafos constituintes do \ac{CORE-SG} são tratados como problemas separados, porque falham por razões diferentes e admitem soluções diferentes. O \ac{k-NNG} é decomponível, mas sua decomposição custa qualidade; a \ac{MST} não é decomponível de forma barata, mas pode ser substituída por uma estimativa. Tratá-los em conjunto obrigaria a aceitar a pior restrição de cada um.\\

Para o \ac{k-NNG}, a construção é reformulada como um conjunto de subproblemas independentes, chamados instâncias. Uma permutação aleatória dos pontos é cortada em $m$ blocos; cada bloco, acrescido de uma amostra de sobreposição extraída dos demais blocos, forma uma instância. Uma instância reúne seus próprios vetores e suas próprias listas de vizinhança, não acessa nenhuma estrutura global durante a execução, e pode portanto ser despachada a qualquer acelerador disponível sem comunicação entre dispositivos. A memória que um dispositivo precisa manter passa a ser determinada pela quantidade de instâncias, independentemente da quantidade de pontos, o que responde à QP1: processar um conjunto de dados maior passa a exigir um valor maior de $m$.\\

O particionamento deixa pares não comparados, e um passo de reparo sobre o grafo mesclado os recupera. Esse passo é o custo dominante da decomposição. A amostra de sobreposição é o parâmetro que responde à QP2: elevá-la faz com que parte dos pares que cruzariam fronteiras seja comparada dentro das instâncias, o que desloca trabalho do reparo, executado na \ac{CPU}, para a fase de instâncias, executada nos dispositivos.\\

Para a \ac{MST}, o trabalho substitui o cálculo exato por uma estimativa, respondendo à QP3. A observação que a justifica é que as únicas arestas que a vizinhança não é capaz de fornecer são aquelas que ligam componentes separados por densidade, e que tais arestas partem necessariamente das bordas desses componentes, pois um ponto interior sempre tem um vizinho mais próximo na direção do percurso. Estimar bordas por meios geométricos é impraticável em alta dimensionalidade, mas elas podem ser estimadas diretamente por uma estatística de \textit{hubness} invertida e normalizada localmente: o grau reverso de um vértice, isto é, quantos outros vértices o listam como vizinho, dividido pelo grau reverso médio de seus próprios vizinhos. Conectar apenas os vértices de menor pontuação nessa estatística é suficiente, e o tamanho desse conjunto é fixado pela dimensionalidade e pela quantidade de arestas do grafo, independentemente do número de pontos.\\

As contribuições deste trabalho são, portanto:\\

\begin{enumerate}
    \item Uma decomposição da construção do \ac{k-NNG} em instâncias independentes, cuja residência por dispositivo é função da quantidade de instâncias, independentemente do tamanho do conjunto de dados, acompanhada da caracterização do que essa decomposição custa em qualidade e de um parâmetro que controla esse custo (Capítulos \ref{cha:methods} e \ref{cha:results});

    \item Uma estimativa da metade de árvore geradora por pontes anti-hub, cujo conjunto abrangido é dimensionado pela dimensionalidade e pela quantidade de arestas, independentemente do número de pontos, com avaliação contra a árvore exata que ela substitui (Seções \ref{sec:bridges} e \ref{sec:res-bridges});

    \item Uma caracterização empírica dos limites do método, incluindo quatro estratégias de redução do passo de reparo que foram medidas e não funcionaram, o que delimita onde o trabalho futuro deve ser dirigido (Seção \ref{sec:res-repair});

    \item Uma quantificação do efeito de dispensar a metade de árvore geradora em grafos conexos, questão levantada teoricamente por \citeonline{core-sg} e, até onde pudemos verificar, não medida anteriormente (Seção \ref{sec:res-connectivity}).
\end{enumerate}

% ---
\section{Relevância\label{sec:significance}}
% ---

A relevância do trabalho tem três planos.\\

No plano imediato, ele amplia a faixa de conjuntos de dados sobre os quais a análise baseada em densidade com varredura de $\mpts$ é executável em aceleradores. Um grafo sobre $5 \times 10^7$ pontos, que em construção não decomposta exigiria dezenas de gigabytes de memória de dispositivo, passa a ser construído com residência por dispositivo da ordem de centenas de megabytes.\\

No plano metodológico, a noção de instância autocontida é independente do \ac{CORE-SG}. Ela se aplica a qualquer construção de grafo de vizinhança por refinamento iterativo, e a propriedade de não exigir comunicação entre dispositivos durante a execução torna o escalonamento direto. A caracterização de quanto essa independência custa, e de como a sobreposição reduz esse custo, transfere-se a outros métodos que enfrentem a mesma escolha.\\

No plano das limitações, o trabalho documenta o que foi medido e não funcionou, e as direções que permanecem abertas.\\

% ---
\section{Delimitação do escopo\label{sec:scope}}
% ---

Este trabalho trata da construção do \ac{CORE-SG} sobre um conjunto de dados fixo e disponível em sua totalidade. A manutenção da estrutura sob chegada contínua de dados, isto é, sobre fluxos, não faz parte do escopo, embora tenha sido objeto de investigação ao longo do projeto. A abstração de instância desenvolvida aqui não exige que um bloco seja novo, e uma instância cujo contexto é extraído de um grafo já existente estende esse grafo em vez de construí-lo; resultados preliminares nessa direção, obtidos antes da delimitação do escopo atual, são apresentados na Seção \ref{sec:fut-streams} como ponto de partida para trabalho futuro, fora do conjunto de contribuições validadas aqui.\\

Também não fazem parte do escopo a remoção de pontos, o tratamento de desvio de conceito, e a execução distribuída entre múltiplos nós. As duas primeiras são consequências naturais do recorte sobre conjuntos fixos; a terceira é discutida na Seção \ref{sec:fut-multinode}, onde se identifica qual passo da construção constitui o obstáculo e quais caminhos permanecem abertos.\\

% ---
\section{Objetivos\label{sec:objectives}}
% ---

O objetivo principal deste trabalho é tornar a construção do \ac{CORE-SG} executável em múltiplos aceleradores, de forma que a memória exigida de cada dispositivo seja determinada por um parâmetro sob controle do usuário, independentemente do tamanho do conjunto de dados, viabilizando a construção de grafos que não caberiam em um único dispositivo.\\

Os objetivos secundários são:
\begin{itemize}
    \item Desenvolver uma decomposição da construção do \ac{k-NNG} cuja unidade de trabalho seja autocontida o bastante para ser distribuída entre dispositivos sem comunicação entre eles, e cuja memória individual seja controlável;

    \item Caracterizar o custo dessa decomposição, tanto em qualidade de vizinhança quanto em tempo de execução, e identificar o parâmetro que permite ao usuário escolher sua posição nesse compromisso;

    \item Desenvolver uma alternativa ao cálculo exato da \ac{MST} sobre o grafo completo, dado que esse cálculo é quadrático e centralizado, e avaliar a qualidade da estimativa resultante contra a árvore que ela substitui;

    \item Avaliar a qualidade das aproximações obtidas contra referências exatas, e não contra rótulos de classe, medindo a razão de peso entre a árvore geradora aproximada e a exata, a concordância entre os agrupamentos planos extraídos de ambas, e a revocação da vizinhança aproximada;

    \item Caracterizar o comportamento do método sob variação de escala, de dimensionalidade, de quantidade de instâncias e de quantidade de dispositivos, e identificar os limites em que suas afirmações deixam de valer.
\end{itemize}

% ---
\section{Organização do texto\label{sec:text-organization}}
% ---

O restante do texto está organizado da seguinte forma. O Capítulo \ref{cha:theory-framework} apresenta o referencial teórico, detalhando o funcionamento do agrupamento baseado em densidade e a construção do \ac{CORE-SG}, formalizando o problema atacado, sintetizando a literatura sobre construção aproximada de grafos de vizinhança, agrupamento distribuído e \textit{hubness}, revisando as abordagens que foram inicialmente consideradas e posteriormente descartadas, e justificando as que foram adotadas. O Capítulo \ref{cha:methods} descreve os materiais e métodos, incluindo a especificação do algoritmo proposto, seu pseudocódigo, e o protocolo experimental com as métricas, conjuntos de dados e linhas de base utilizados. O Capítulo \ref{cha:results} apresenta os resultados, organizados na ordem das questões de pesquisa. O Capítulo \ref{cha:discussion} discute o significado e o alcance desses resultados, delimita o escopo das afirmações feitas, apresenta os resultados preliminares sobre fluxos de dados, e enumera os desdobramentos de pesquisa que o trabalho deixa em aberto.\\
